%% kapa: exam=1 round=19 year=2008 period=2 form=A subject=economics subject_ko=경제원론 seq=1 number=1
%% source: Q-Net 감정평가사 기출문제 (cst00309) artlSeq=1001962 file=extracted/1008/1001962/2008년 제19회 감정평가사 1차시험 기출문제/2교시.hwp
%% conversion: kordoc 4.24.51 -> tools/gen1.py (auto); review=proofread-sonnet; flags=none
\begin{question}{1}
놀부는 흥부와 1시간 동안 비디오를 보는 대신에 아르바이트를 해서 5,000원을 벌었다. 놀부가 아르바이트를 하는 것의 기회비용은?
\begin{choices}
  \item 아르바이트로 번 5,000원이다.
  \item 비디오를 보았다면 얻을 수 있었던 즐거움이다.
  \item 5,000원에서 비디오를 보았다면 얻을 수 있었던 즐거움을 뺀 값이다.
  \item 비디오를 보았다면 얻을 수 있었던 즐거움에서 5,000원을 뺀 값이다.
  \item 비디오를 보았다면 5,000원보다 더 작은 즐거움을 얻었을 것이므로 기회비용은 0원이다.
\end{choices}
\end{question}
\answer{2}
